В настоящем обзоре исследован комптоновский процесс в присутствии замагниченной зарядово-симметричной плазмы $\gamma e \to \gamma e$ как с учетом конечной ширины поглощения электрона, так и в случае узкого резонансного пика. 

Представлен детальный сравнительный анализ сечения комптоновского рассеяния, полученного с использованием дельта-функционального приближения~\cite{Rumyantsev:2017} для относительно слабого магнитного поля $B\lesssim B_e$, с известными в~литературе~\cite{Harding:1986,Mushtukov:2016,SchwarmD:2017} результатами, учитывающими конечную ширину резонансного пика. Показано, что дельта-функциональное приближение достаточно хорошо описывает сечение комптоновского рассеяния в области резонансного пика, и точность приближения тем лучше, чем ниже температура и больше значение напряженности магнитного поля. Выражение для сечения в пределе узкого резонансного пика имеет простой вид, удобный для численного анализа, что позволило рассчитать сечение для случая, когда начальный электрон может занимать произвольный уровень Ландау. Численный анализ показал, что учет ненулевых уровней Ландау начального электрона приводит к уширению резонансных пиков и к завышенным значениям сечения  вблизи резонанса, что становится особенно существенным для относительно высоких температур $T\gtrsim 50$ кэВ.

В пределе сильного магнитного поля $B>B_e$ и высоких температур \linebreak \mbox{$T=1$ МэВ} вычислен коэффициент поглощения фотона в комптоновском процессе для двух кинематически разрешенных каналов \mbox{$\gamma^{(1)} e\to \gamma^{(1)} e$}\linebreak и $\gamma^{(1)} e\to \gamma^{(2)} e$ . Полученные численные результаты позволили установить области энергий фотона, при которых применим нерезонансный предел~\cite{Chistyakov:2009}:\linebreak $\omega\lesssim 3$ МэВ при $B=200B_e$ и $\omega\lesssim 0.3$ МэВ при $B=20B_e$. Дельта-функциональное приближение при данных параметрах среды хорошо описывает только первый резонансный пик.  Как известно, в замагниченной среде возможно затухание фотона как электромагнитной волны, поэтому представляет интерес исследовать возможность формирования комптоновского процесса, что является целью следующей главы.

Исследован процесс распространения квантованной электромагнитной волны в сильно замагниченной зарядово-симметричной плазме. С учетом изменения дисперсионных свойств фотона в магнитном поле и плазме было установлено, что, аналогично случаю чистого магнитного поля, процесс затухания фотона
в замагниченной плазме имеет неэкспоненциальный характер. 

Для характерного отрезка времени $\sim [W^{{(\lambda)}}_{abs}]^{-1}$ была использована аппроксимация экспоненциально затухающими колебаниями. В этом случае было
показано, что коэффициент поглощения фотона в околопороговой области меньше по сравнению с известными в литературе результатами. Также, следуя данной аппроксимации, были построены дисперсионные кривые, которые и для моды~1, и для моды 2 близки к вакуумным кривым, за исключением околопороговых областей. Полученные результаты согласуются с выводами работы~\cite{Shabad:1988}.

Кроме того, был выполнен анализ возможности формирования комптоновского процесса в условиях нестабильности фотона за счет процесса поглощения электроном $e^\pm\gamma\to e^\pm$ и распада на электрон-позитронную пару $\gamma\to e^+e^-$. Несмотря на малый фактор $\alpha$, комптоновский процесс преобладает над процессами распада и поглощения в области низких энергий фотона ($\omega<3$ МэВ при $T=1$ МэВ и $B =200B_e$). С другой стороны, даже с учетом резонанса на виртуальном электроне фотоны как моды 1, так и моды~2 в пределе сильного магнитного поля $B=200 B_e$ и при температуре $T=1$ МэВ эффективно затухают в резонансной области.

Для определения характера затухания фотона необходимо выполнить более детальный анализ временной зависимости амплитуды $F_A(t)$ колебаний, входящей в уравнение~(\ref{eq:Fm}).